\documentclass[11pt,a4paper]{article}

\usepackage{kotex}
\usepackage{fontspec}
\setmainfont{Times New Roman}
\setmainhangulfont{AppleMyungjo}
\setsanshangulfont{Apple SD Gothic Neo}

\usepackage[a4paper,top=18mm,bottom=18mm,left=20mm,right=20mm]{geometry}
\usepackage{fancyhdr}
\setlength{\headheight}{24pt}
\usepackage{enumitem}
\usepackage{mdframed}
\usepackage{array}
\usepackage{booktabs}

\setlength{\parindent}{1em}
\linespread{1.18}

\pagestyle{fancy}
\fancyhf{}
\renewcommand{\headrulewidth}{0.6pt}
\fancyhead[L]{\sffamily\bfseries 2026 고1 영어 변형 — 지문 19 [정답 및 해설]}
\fancyhead[R]{\sffamily Daniel \& 기차 지연}
\renewcommand{\footrulewidth}{0pt}
\fancyfoot[C]{\framebox[16mm][c]{\thepage}}

\newcommand{\qno}[1]{\par\vspace{10pt}\noindent{\sffamily\bfseries #1.}\ }
\newcommand{\instr}[1]{{\sffamily #1}\par\vspace{3pt}}
\newlist{choices}{enumerate}{1}
\setlist[choices]{label=,leftmargin=1.5em,itemsep=2pt,topsep=3pt,parsep=0pt}
\newmdenv[linewidth=0.6pt,roundcorner=3pt,innertopmargin=7pt,innerbottommargin=7pt,
          innerleftmargin=9pt,innerrightmargin=9pt,skipabove=5pt,skipbelow=5pt]{pbox}

\begin{document}

\begin{center}
{\fontsize{15}{17}\selectfont\sffamily\bfseries [지문 19] 정답 및 해설}
\end{center}
\vspace{8pt}

%% 정답 표
\begin{center}
\begin{tabular}{>{\centering\arraybackslash}p{2cm}
                >{\centering\arraybackslash}p{3.5cm}
                >{\centering\arraybackslash}p{2cm}
                >{\centering\arraybackslash}p{2cm}}
\toprule
\bfseries 문항 & \bfseries 유형 & \bfseries 정답 & \bfseries 배점 \\
\midrule
1 & 빈칸추론 & ② & 2점 \\
2 & 함의추론 & ③ & 2점 \\
3 & 서술형 (해석) & 모범답안 참조 & 3점 \\
4 & 심경 변화 & ② & 2점 \\
\bottomrule
\end{tabular}
\end{center}

\vspace{12pt}

%% 문항 1 해설
\qno{1}\instr{[빈칸추론] 정답: ②  apprehensive}

\noindent\textbf{정답 근거}: Daniel은 버스가 갑자기 고장나는 바람에 기차를 놓칠까 봐 걱정하게 된 상황이다.
빈칸 뒤에 ``about missing his train''이 이어지므로, 불안·염려를 나타내는 형용사가 필요하다.
\textit{apprehensive}(불안한, 염려하는)가 문맥에 가장 잘 맞는다.

\vspace{4pt}
\noindent\textbf{오답 소거}:
\begin{choices}
\item ① \textit{indifferent}(무관심한) — 기차를 놓칠 위기에 무관심할 이유가 없다.
\item ③ \textit{elated}(기분이 들뜬) — 기쁜 감정으로, 불안한 상황과 정반대.
\item ④ \textit{relieved}(안도한) — 안도는 역에 도착한 뒤의 결말 감정이며, 이 시점에서는 아직 불안 상태다.
\item ⑤ \textit{triumphant}(의기양양한) — 성취감을 나타내며, 위기 상황에서 쓰기 어렵다.
\end{choices}

\vspace{10pt}

%% 문항 2 해설
\qno{2}\instr{[함의추론] 정답: ③}

\noindent\textbf{정답 근거}: Daniel은 역에 도착해 기차가 10분 지연됐다는 것, 즉 기차가 자신을 두고 출발하지 않았음을 확인한 뒤 비로소 긴 안도의 숨을 내쉰다.
따라서 밑줄 친 표현은 ``기차가 자기를 두고 떠나지 않았다는 사실을 알고 긴장이 풀렸다''는 의미이다.

\vspace{4pt}
\noindent\textbf{오답 소거}:
\begin{choices}
\item ① 역까지 달려온 탓에 지쳐서 — 본문에는 달려갔다는 묘사가 없다.
\item ② 택시를 탄 빠른 결단이 자랑스러워서 — 자긍심이 아닌 안도감이 핵심이다.
\item ④ 기차가 정시에 오지 않아 실망해서 — 지연은 오히려 Daniel에게 유리한 소식이었다.
\item ⑤ 기계 결함이 수리됐기 때문에 — 본문은 기계 결함 수리 여부를 언급하지 않는다.
\end{choices}

\vspace{10pt}

%% 문항 3 해설
\qno{3}\instr{[서술형] 정답 (모범답안)}

\begin{pbox}
\textbf{대상 문장}: Without hesitation, he got off the bus and took a taxi to the station.

\vspace{6pt}
\textbf{모범답안}: 주저하지 않고, 그는 버스에서 내려 역까지 택시를 탔다.
\end{pbox}

\vspace{6pt}
\noindent\textbf{채점 포인트} (총 3점):
\begin{choices}
\item \textbullet\ ``Without hesitation'' $\rightarrow$ ``주저(망설임) 없이 / 거리낌 없이'' \hfill (1점)
\item \textbullet\ ``got off the bus'' $\rightarrow$ ``버스에서 내렸다'' \hfill (1점)
\item \textbullet\ ``took a taxi to the station'' $\rightarrow$ ``역까지(역으로) 택시를 탔다'' \hfill (1점)
\end{choices}

\vspace{4pt}
\noindent\textbf{감점 기준}: 세 항목 중 하나라도 의미가 왜곡되거나 누락되면 해당 항목 점수 미부여.
자연스럽지 않은 직역(예: ``주저함 없이 그는 버스로부터 내렸다'')은 0.5점 감점 가능.

\vspace{8pt}
\noindent{\sffamily\bfseries 4. 심경 변화 - 정답: ②}
\par Daniel은 버스 고장 후 기차를 놓칠까 봐 불안해했지만, 기차가 10분 지연되었다는 사실을 알고 안도한다. 따라서 \textit{anxious $\rightarrow$ relieved}가 가장 적절하다.

\end{document}
